% ---
% id: transmitter/ass:seq
% kind: ass
% title: "Sequential round structure"
% summary: >-
%   Each round runs the delivery of its quantum systems, then every party's
%   measurement, then the round's public messages, then sifting; each party
%   measures in round order, post-processing follows all rounds and every
%   message is public. It is the timing premise of the security analysis.
% prominence: standard
% status: accepted
% uses: []
% assumes: []
% refs:
%   - "tupkary2026decoy, Def. 8.1, Def. 8.2, Lemma 8.2"
%   - "tupkary2026decoy, Sec. 8.1.1, Protocol 2"
% provenance: newly-derived
% short_id: A-SEQ
% verification: never
% category: protocol-structure
% affects: [secrecy]
% user_statement: >-
%   In every round the quantum systems come first: where your source is
%   trusted, it prepares and sends that round's signal, and otherwise the
%   systems of the round are delivered to the measuring stations. Then every
%   party that measures completes its measurement, and only then are that
%   round's public messages sent, in the declared order. Sifting is then
%   determined by the announcements, and no operation on a quantum system of a round
%   depends on a public message of that round. All classical post-processing
%   (for example error correction, verification and privacy amplification)
%   runs after all n rounds, and every classical message is public. A message
%   sent before a measurement of the same round has finished (for example an
%   early basis announcement) violates this. Each party that measures
%   completes its measurement of each round before it completes that of the
%   next round.
% lean:
% ---

With one measuring party this is the structure of
\cite[Def.~8.1, Def.~8.2]{tupkary2026decoy}, whose state evolution is its
Lemma~8.2, and it is the hypothesis under which this module applies the
imported source-replacement lemma, whose own statement places no condition on
the sequential structure of the operations after the preparation. The cited definitions are that case and
no more: they have one measuring party and one announcement block per round,
and the reading for several measuring parties is this contract's own. Of the
messages, the assumption fixes only that all of them follow all measurements
of the round; a message may depend on earlier messages of the same round,
which is a property of the declared messages and not of this assumption. The
clauses have two consequences worth naming. Where one party
measures, as in a prepare-and-measure protocol, every measurement of the round
is that party's. Where two parties measure systems delivered by a source
neither of them prepared, neither party's message precedes the other party's
measurement. Clause (v) is the measurement part of the timing condition of
the published Protocol~2 (Bob's measurement times increasing there), read for
each measuring party \cite[Sec.~8.1.1, Protocol~2]{tupkary2026decoy}; the
monotone order of preparations and of announcements is not adopted, and the
analysis does not need it; with (i) it puts every
measurement of a round before every public message of a later round.

\begin{assumption}\label{ass:seq}
The physical protocol has the sequential structure the declared protocol
describes:
(i) each round runs, in this order, state preparation and transmission; the
measurements of the round, by every party that measures in it; the public
messages of the round, in their declared order, \emph{each of them sent after
every measurement of the round is complete}, whichever party sends it and
whichever party measured; and announcement-determined sifting;
(ii) the classical post-processing is applied after the $n$ rounds;
(iii) every message is public;
(iv) no operation on a quantum system of round $j$ depends on a public
message of round $j$;
(v) every party that measures makes its measurements in round order: its
measurement of round $j$ is complete before its measurement of round $j+1$
is.
\end{assumption}
